\scanpage{277}
et alors $\uuline{\operatorname{N\acute{e}r}}_{X/S}$ est représenté par $\uline{\operatorname{N\acute{e}r}}_{X/S}$

\uline{Corollaire 1} S'il en est ainsi, $\uline{\operatorname{N\acute{e}r}}^{\deleted{\unsure{\tau}}}_{X/S}$ est \deleted{\unsure{représentable}} séparé ssi $\uline{\operatorname{Pic}}^{\tau}_{X/S}$ est fermé dans $\uline{\operatorname{Pic}}_{X/S}$. En tous cas, $\uline{\operatorname{N\acute{e}r}}_{X/S}$ est \uline{non ramifié} sur $S$.

\uline{Corollaire 2} Supposons que $\uline{\operatorname{Pic}}^{\tau}_{X/S}$ soit \uline{propre} \inserted{et \uline{plat}} sur $S$, et que $\uline{\operatorname{Pic}}^{\deleted{\unsure{\tau}}}_{X/S}$ soit réunion \inserted{croissante} d'ouverts \inserted{$U_i$} stables par $\uline{\operatorname{Pic}}^{\tau}_{X/S}$ et quasi-projectifs sur $S$. [cette dernière condition est conséquence de $\uline{\operatorname{Pic}}^{\tau}_{X/S}$ de type fini sur $S$, dans le cas où $X$ est projectif, plat,\inserted{\deleted{\unsure{plat et}} $\unsure{f_*(\mathcal O_X)=\mathcal O_S}$ universellement} et primaire sur $S$]. Alors $\uline{\operatorname{N\acute{e}r}}_{X/S}$ existe et est séparé non ramifié sur $S$.
\margin{La première condition est vérifiée si les fibres de $X/S$ ont des \unsure{Nér géo} sans torsion et des $\operatorname{Pic}^0$ propres\ldots}

Il semble que dès que $\uline{\operatorname{Pic}}^{\tau}_{X/S}$ est \uline{plat} \deleted{\unsure{sur}} sur $S$\inserted{et fermé dans $\uline{\operatorname{Pic}}_{X/S}$}, $\uline{\operatorname{N\acute{e}r}}_{X/S}$ doive exister : c'est lié à la question de la caractérisation des foncteurs représentés par un objet non ramifié sur $S$\ldots. Au concours ! [A vrai dire, même l'existence de $\uline{\operatorname{Pic}}^{\deleted{\unsure{\tau}}}_{X/S}$ ne semble pas indispensable !]
\margin{?}
\figure{\begin{tikzpicture}[x=1cm,y=1cm]
\draw (0,7)--(0.02,6.15); \draw (0.12,5)--(0.13,2.25); \draw (0.08,1.25)--(0,0.65)--(0.12,-0.45);
\draw (2.5,3.12)..controls (2.1,3.5) and (4.4,3.4)..(4.85,3.35)..controls (5.1,3.1) and (4.8,3.0)..(4.6,3.0)--(2.6,2.94)--(2.55,2.65);
\draw (1.1,1.4)..controls (0.7,1.7) and (2.5,1.75)..(2.8,1.5)--(2.8,1.25)--(1.2,1.16)--(1.15,0.98);
\end{tikzpicture}}
\mspage{276}
